\section{Joint restriction types and multiplication on two sides}
\label{l:two-sided-section}

For a fixed collection of disjoint block subgroups, we can retain the
number of distinct joint restriction types rather than summing all
small factor degrees.  This gives an operator estimate by orbit spans.
It also gives a product estimate when one measure is capped on left
cosets and another is capped on right cosets.  These two conclusions
use the same type count, but different Fourier estimates.

Let $M_1,\ldots,M_b$ be disjoint nonempty subsets of $\{1,\ldots,n\}$,
and put $H_i=\Sym(M_i)$, acting trivially on the complement.  On an
irreducible $V_\lambda$ of $S_n$, restriction to
$H=H_1\times\cdots\times H_b$ has the form
\begin{equation}\label{l:joint-decomposition}
 V_\lambda|_H
   =\bigoplus_{\boldsymbol\tau}
       (V_{\tau_1}\otimes\cdots\otimes V_{\tau_b})
                       \otimes\mathcal A_{\boldsymbol\tau}.
\end{equation}
The sum ranges over distinct occurring joint types; the nonzero
space $\mathcal A_{\boldsymbol\tau}$ records their arbitrary
multiplicities.  Write $K_\lambda$ for the number of terms in this
sum and $Q_{\boldsymbol\tau}$ for their orthogonal projections.

\begin{lemma}[Counting joint types]
\label{l:joint-types}
Put $k=n-\max(\lambda_1,\lambda'_1)$ and
$M(k)=\sum_{j=0}^kp(j)$, where $p(0)=1$ and $p(j)$ is the
partition number.  Then
\begin{equation}\label{l:joint-types-bound}
 K_\lambda\le M(k)^b.
\end{equation}
For every fixed $b$ and every fixed $\eta>0$,
\begin{equation}\label{l:joint-types-summability}
 \sum_{\lambda\vdash n:\,D_\lambda>1}
         K_\lambda^2D_\lambda^{-\eta}\longrightarrow0.
\end{equation}
The statement is uniform over the disjoint sets $M_i$ and does not
require them to have equal sizes.
\end{lemma}
\begin{proof}
By branching, a type of $H_i$ has shape $\nu\subseteq\lambda$
and size $|M_i|$.  Orient $\lambda$ so that a longest line is a
row.  Every such $\nu$ has at most $k$ boxes below its first row.
For a fixed lower size $j$, the lower partition determines the first
row because $|\nu|=|M_i|$ is fixed.  There are therefore at most
$M(k)$ possibilities for each factor.  If the longest line was a
column, transpose all diagrams.  This proves
\eqref{l:joint-types-bound}; absent tuples only reduce the count.

There are at most $2p(k)$ ambient diagrams with this deficit.  The
partition bound $p(j)\le e^{3\sqrt j}$ gives, for a constant $C_b$,
\[
 2p(k)M(k)^{2b}\le e^{C_b\sqrt k}\qquad(k\ge1).
\]
By Theorem~\ref{l:degree-all-powers},
$D_\lambda\ge c_1e^{c_2k}$ and, at each fixed positive $k$,
the dimensions tend uniformly to infinity.  Thus the contribution
at deficit $k$ is bounded by
$c_1^{-\eta}e^{C_b\sqrt k-\eta c_2k}$, a summable sequence,
and tends to zero for fixed $k$.  Dominated convergence proves
\eqref{l:joint-types-summability}.  Multiplicities do not affect
either counting step.
\end{proof}

\subsection{An operator estimate retaining the joint-type count}

\begin{theorem}[Joint-type orbit transfer]
\label{l:joint-type-operator}
Let $f\ge0$ be a mass function on $S_n$, of total mass $z\le1$,
such that $f(gH_i)\le B/[S_n:H_i]$ for every $g,i$.  With
$K_\lambda$ as in \eqref{l:joint-decomposition},
\begin{equation}\label{l:joint-type-operator-bound}
 \|\widehat f(\lambda)\|_{\op}
       \le\sqrt{BzK_\lambda}\,D_\lambda^{-1/2+1/b}.
\end{equation}
\end{theorem}
\begin{proof}
First let $v$ lie in a single $H_i$-isotypic space of type $\tau_i$.
The orbit span
$W=\Span\{\rho_\lambda(h)v:h\in H_i\}$ has dimension at most
$d_{\tau_i}^2$.  To see this in the presence of multiplicity, write
the full $H_i$-isotypic space as $V_{\tau_i}\otimes\mathcal A$;
the orbit lies in the image of the linear map
$A\mapsto(A\otimes I)v$ from $\operatorname{End}(V_{\tau_i})$.

For a test vector $w$,
\[
 |\langle w,\rho_\lambda(g)v\rangle|
    \le\|v\|\,\|P_{\rho_\lambda(g)W}w\|.
\]
The squared projected norm is constant on each $gH_i$, and its
uniform mean is $(\dim W/D_\lambda)\|w\|^2$: the average of
the conjugate projections is scalar by Schur's lemma, and its trace
is $\dim W$.  Cauchy--Schwarz for the mass function $f$, followed
by its coset cap, gives
\[
 |\langle w,\widehat f(\lambda)v\rangle|^2
 \le zB\frac{\dim W}{D_\lambda}\|v\|^2\|w\|^2
 \le\frac{zB d_{\tau_i}^2}{D_\lambda}\|v\|^2\|w\|^2.
\]
Thus this single-type bound holds before any joint decomposition.
In a joint type of \eqref{l:joint-decomposition}, the factor degrees
satisfy $\prod_i d_{\tau_i}\le D_\lambda$, since the entire tensor
representation occurs at least once.  Choose $i$ with
$d_{\tau_i}\le D_\lambda^{1/b}$ and apply the preceding bound.
For an arbitrary $v$, decompose it into its $K_\lambda$ orthogonal
joint components and sum their bounds.  Their norms have sum at
most $\sqrt{K_\lambda}\|v\|$, proving
\eqref{l:joint-type-operator-bound}.
\end{proof}

The following two specializations keep two useful ways of comparing
the joint-type count with the ambient dimension.  The four-block
argument uses a uniform linear lower bound in the broad-diagram
range.  The eight-block argument uses a different cutoff for the
longest line and a direct hook estimate below that cutoff.

\begin{corollary}[Four-block transfer]
\label{l:four-block-reservoir}
Partition $n$ points into four sets of size $n/4$, and let $H_i$
permute those sets.  For all sufficiently large such $n$, every
subprobability $f$ with $f(gH_i)\le2/[S_n:H_i]$ satisfies
\[
 \|\widehat f(\lambda)\|_{\op}\le D_\lambda^{-1/8}
       \qquad(D_\lambda>1).
\]
\end{corollary}
\begin{proof}
Theorem~\ref{l:joint-type-operator} gives
$\|\widehat f(\lambda)\|_{\op}
 \le\sqrt{2K_\lambda}D_\lambda^{-1/4}$, and
Lemma~\ref{l:joint-types} gives $\log(2K_\lambda)=O(\sqrt k)$.
We verify directly that $\log D_\lambda/\sqrt k\to\infty$
uniformly over $k\ge1$.

Orient a longest line as the first row, of length $a=n-k$.
For $1\le k\le n/3$, the tableau construction of
Lemma~\ref{l:tableau-constructions} gives
\[
 D_\lambda\ge\binom{n-k}k\ge((n-k)/k)^k.
\]
For $k$ above a fixed cutoff, divide the logarithm by $\sqrt k$
and use $(n-k)/k\ge2$; for the finitely many smaller positive
$k$, the same expression tends to infinity with $n$.
For $k>n/3$ and $a\le n/10$, the hook formula gives
$D_\lambda\ge(n/(2ea))^n\ge(5/e)^n$.
If instead $a>n/10$, retain the first row and
$r=\lfloor a/2\rfloor$ lower boxes; there are enough lower boxes
because $k>n/3>a/2$.  Their tableaux extend to the full diagram,
and the same construction gives
$D_\lambda\ge\binom ar\ge2^r\ge e^{cn}$ for an absolute
$c>0$ and all large $n$.  This proves the asserted uniform divergence.
Consequently $2K_\lambda\le D_\lambda^{1/4}$ eventually,
which yields the displayed operator bound.
\end{proof}

\begin{corollary}[Eight-block transfer]
\label{l:eight-block-probe}
Partition $n$ points into eight sets of size $n/8$, and let $H_i$
permute those sets.  For all sufficiently large such $n$, every
subprobability $f$ with $f(gH_i)\le2/[S_n:H_i]$ satisfies
\[
 \|\widehat f(\lambda)\|_{\op}\le D_\lambda^{-1/4}
       \qquad(D_\lambda>1).
\]
\end{corollary}
\begin{proof}
Here the joint-type estimate is
$\sqrt{2K_\lambda}D_\lambda^{-3/8}$.
For $k\le n/3$, use $K_\lambda\le M(k)^8$ and
$D_\lambda\ge\binom{n-k}k$ exactly as above to get
$\log(2K_\lambda)/\log D_\lambda\to0$ uniformly.
For $k>n/3$, the equal block sizes give the coarser count
$K_\lambda\le p(n/8)^8$, whose logarithm is $O(\sqrt n)$.
Let $a$ again be the longest line.  If $a\ge n^{3/4}$, retain
that line and $r=\lfloor a/3\rfloor$ lower boxes, possible because
$k>n/3>a/3$.  Tableau extension gives
$D_\lambda\ge\binom ar\ge3^r$, so
$\log D_\lambda\ge c n^{3/4}$.  If $a<n^{3/4}$, all hooks
are at most $2n^{3/4}$, and
\[
 D_\lambda\ge\frac{n!}{(2n^{3/4})^n}
       \ge\left(\frac{n^{1/4}}{2e}\right)^n.
\]
Thus $2K_\lambda\le D_\lambda^{1/4}$ uniformly for all large
$n$, and the result follows.
\end{proof}

These are exact specializations of the joint-type orbit argument,
including its multiplicity treatment.  For reference, the four-block
proof also gives $\sum_{D_\lambda>1}D_\lambda^{-4}=o(1)$ directly:
at $k\le n/3$ use the summable majorant
$2e^{3\sqrt k}2^{-4k}$ and fixed-$k$ divergence; at $k>n/3$
use the linear-in-$n$ logarithmic dimension bound against $p(n)$.
The eight-block proof gives the inverse-third-power statement by
the same summation, with majorant $2e^{3\sqrt k}2^{-3k}$ and its
$n^{3/4}$ lower bound in the remaining range.  With mass and sign
input, these reproduce respectively the 24-factor and 10-factor
Fourier completions through Proposition~\ref{l:amplification}:
the nonlinear Plancherel exponents are respectively $-4$ and $-3$.
Alternatively, a final independent probability law with zero sign
mean removes the sign coefficient exactly.  Its Fourier matrices
have operator norm at most one, so this extra convolution preserves
both nonlinear estimates.  This alternative does not require small
sign mean in the law supplying the coset caps.

\subsection{A product with caps in opposite orientations}

For a mass function $\alpha$, a cap on $H_i g$ controls left
multiplication of its Fourier matrix by an $H_i$-isotypic projector.
A cap on $gH_i$ controls right multiplication.  The product below
places these two controlled sides together.  There is no assertion
that either cap implies the other.

\begin{theorem}[Two-sided Fourier multiplication]
\label{l:two-sided}
Let $H_1,\ldots,H_b$ be the disjoint symmetric block subgroups
above.  Let $\alpha,\beta\ge0$ have masses $a,c\le1$ and satisfy
\begin{equation}\label{l:two-sided-caps}
 \alpha(H_i g)\le\frac{B_L}{[S_n:H_i]},\qquad
 \beta(gH_i)\le\frac{B_R}{[S_n:H_i]}
       \quad(g\in S_n,\ 1\le i\le b).
\end{equation}
Then every irreducible $\lambda\vdash n$ satisfies
\begin{equation}\label{l:two-sided-bound}
 D_\lambda\|\widehat\beta(\lambda)
                     \widehat\alpha(\lambda)\|_{\HS}^2
 \le B_LB_Rac\,K_\lambda^2D_\lambda^{-1+4/b}.
\end{equation}
In particular, for eight blocks and $B_L=B_R=2$, the bound is
$4K_\lambda^2D_\lambda^{-1/2}$.
\end{theorem}
\begin{proof}
Use densities $A=|S_n|\alpha$ and $B=|S_n|\beta$ with respect
to uniform probability on $G=S_n$.  Write $\int_G$ and
$\int_{H_i}$ for normalized counting averages.  The two caps read
\begin{equation}\label{l:two-sided-density-caps}
 \int_{H_i}A(h^{-1}g)\,dh\le B_L,
 \qquad
 \int_{H_i}B(gh^{-1})\,dh\le B_R.
\end{equation}
Indeed these averages are respectively
$[G:H_i]\alpha(H_i g)$ and $[G:H_i]\beta(gH_i)$.
The density transform $\int_G A(g)\rho_\lambda(g)\,dg$ equals
the mass transform $\widehat\alpha(\lambda)$, and similarly for
$B$.  These conventions give
$\widehat{\beta*\alpha}=\widehat\beta\widehat\alpha$.

For a type $\tau$ of $H_i$ of degree $d_\tau$ and character
$\chi_\tau$, its full isotypic projector is
\[
 P_{i,\tau}=d_\tau\int_{H_i}
              \overline{\chi_\tau(h)}\rho_\lambda(h)\,dh.
\]
Centrality makes this scalar on each irreducible carrier, and
character orthogonality identifies the scalar as one on type $\tau$
and zero otherwise.  It acts as the identity on all multiplicity
copies of $\tau$.
Consider the density
\[
 C(g)=d_\tau\int_{H_i}\overline{\chi_\tau(h)}A(h^{-1}g)\,dh.
\]
Weighted Cauchy--Schwarz and nonnegativity give
\[
 |C(g)|^2\le d_\tau^2
  \left(\int_{H_i}A(h^{-1}g)\,dh\right)
  \left(\int_{H_i}|\chi_\tau(h)|^2A(h^{-1}g)\,dh\right).
\]
The first factor is at most $B_L$.  The average over $G$ of the
second factor is $a\int_{H_i}|\chi_\tau|^2=a$.
The Fourier matrix of $C$ is
$P_{i,\tau}\widehat\alpha(\lambda)$.
Plancherel, retaining the $\lambda$ term, therefore yields
\begin{equation}\label{l:two-sided-isotypic}
 D_\lambda\|P_{i,\tau}\widehat\alpha(\lambda)\|_{\HS}^2
       \le B_La d_\tau^2,
 \qquad
 D_\lambda\|\widehat\beta(\lambda)P_{i,\tau}\|_{\HS}^2
       \le B_Rc d_\tau^2.
\end{equation}
For the second inequality use
$d_\tau\int_{H_i}\overline{\chi_\tau(h)}B(gh^{-1})\,dh$,
whose transform is $\widehat\beta(\lambda)P_{i,\tau}$, and the
second cap in \eqref{l:two-sided-density-caps}.

The projectors for the different disjoint factors commute, and
$Q_{\boldsymbol\tau}=\prod_iP_{i,\tau_i}$ are the orthogonal
joint projections in \eqref{l:joint-decomposition}.  On a nonzero
joint space, $\prod_i d_{\tau_i}\le D_\lambda$, so some
$d_{\tau_i}\le D_\lambda^{1/b}$.
Since $Q_{\boldsymbol\tau}\le P_{i,\tau_i}$, the two bounds
in \eqref{l:two-sided-isotypic} imply
\[
 \begin{split}
 \|\widehat\beta Q_{\boldsymbol\tau}\widehat\alpha\|_{\HS}
 &\le\|\widehat\beta Q_{\boldsymbol\tau}\|_{\HS}
             \|Q_{\boldsymbol\tau}\widehat\alpha\|_{\HS}\\
 &\le\sqrt{B_LB_Rac}\,D_\lambda^{-1+2/b}.
 \end{split}
\]
Sum over the $K_\lambda$ joint types before squaring, and then
multiply by $D_\lambda$.  This gives
\eqref{l:two-sided-bound}.  The factor $K_\lambda^2$ is retained;
no multiplicity-free restriction has been assumed.
\end{proof}

\begin{corollary}[Two-factor completion]
\label{l:two-sided-completion}
Fix $b>4$ and finite constants $B_L,B_R$.  Suppose probabilities
$\mu_n,\nu_n$ on $S_n$ dominate subprobabilities
$\alpha_n,\beta_n$, respectively, whose masses tend to one and
which satisfy \eqref{l:two-sided-caps} for $b$ disjoint nonempty
blocks.  If
$\widehat\nu_n(\sgn)\widehat\mu_n(\sgn)\to0$, then
\[
 \|\nu_n*\mu_n-U_{S_n}\|_{\TV}\longrightarrow0.
\]
In particular this applies to a single law with small sign mean
when separate subprobabilities supply its two cap orientations.
\end{corollary}
\begin{proof}
Use Lemma~\ref{l:joint-types} with $\eta=1-4/b>0$ in
Theorem~\ref{l:two-sided}.  It makes the sum of the nonlinear
Plancherel terms for $\beta_n*\alpha_n$ tend to zero.
Its trivial coefficient is the product of the two masses, tending
to one.  Domination bounds the error in each sign coefficient by
the corresponding lost mass, so its sign coefficient also tends
to zero.  Lemma~\ref{l:plancherel} now gives convergence to uniform
in $\ell^1$.  Finally,
$\beta_n*\alpha_n\le\nu_n*\mu_n$ and the difference has mass
one minus the product of the retained masses.  Adding that mass
proves the conclusion.
\end{proof}
